\markdownRendererDocumentBegin
\markdownRendererSectionBegin
\markdownRendererHeadingOne{COMP0082 Coursework Report — Protein Enzyme Classification}\markdownRendererInterblockSeparator
{}\markdownRendererSectionBegin
\markdownRendererHeadingTwo{1. Introduction}\markdownRendererInterblockSeparator
{}Enzymes are biological catalysts that underpin virtually every metabolic process in living\markdownRendererSoftLineBreak
{}organisms. The Enzyme Commission (EC) numbering system classifies enzymes into six top-level\markdownRendererSoftLineBreak
{}classes—Oxidoreductases (EC 1), Transferases (EC 2), Hydrolases (EC 3), Lyases (EC 4),\markdownRendererSoftLineBreak
{}Isomerases (EC 5), and Ligases (EC 6)—according to the reaction they catalyse. Computational\markdownRendererSoftLineBreak
{}prediction of enzyme class from amino acid sequence alone is of substantial practical\markdownRendererSoftLineBreak
{}importance: it accelerates the functional annotation of the millions of uncharacterised\markdownRendererSoftLineBreak
{}proteins generated by high-throughput sequencing, and it supports drug target identification,\markdownRendererSoftLineBreak
{}metabolic engineering, and synthetic biology. As genome sequencing throughput continues to\markdownRendererSoftLineBreak
{}outpace experimental characterisation, automated enzyme classification becomes increasingly\markdownRendererSoftLineBreak
{}essential for bridging the sequence–function gap.\markdownRendererParagraphSeparator
{}This project addresses the problem as a seven-class supervised classification task, where\markdownRendererSoftLineBreak
{}class 0 represents "not an enzyme" and classes 1–6 correspond to the six EC classes. The\markdownRendererSoftLineBreak
{}dataset comprises 39 764 non-homologous protein sequences with a heavily skewed distribution:\markdownRendererSoftLineBreak
{}class 0 constitutes 81.5 \markdownRendererPercentSign{} of the data, while the smallest class (Ligase, EC 6) accounts\markdownRendererSoftLineBreak
{}for only 0.7 \markdownRendererPercentSign{}. This extreme class imbalance—spanning two orders of magnitude between the\markdownRendererSoftLineBreak
{}largest and smallest classes—is the central challenge of the project and informs every design\markdownRendererSoftLineBreak
{}decision, from loss weighting and oversampling strategies to the choice of evaluation metrics.\markdownRendererParagraphSeparator
{}The approach progresses systematically from handcrafted sequence-based features through protein\markdownRendererSoftLineBreak
{}language model (PLM) embeddings to an ensemble that combines the complementary strengths of\markdownRendererSoftLineBreak
{}gradient-boosted trees and end-to-end fine-tuned transformers. All experiments use stratified\markdownRendererSoftLineBreak
{}five-fold cross-validation with strict data-leakage prevention, and all results are reported\markdownRendererSoftLineBreak
{}with standard deviations across folds to quantify variability.\markdownRendererInterblockSeparator
{}
\markdownRendererSectionEnd \markdownRendererSectionBegin
\markdownRendererHeadingTwo{2. Methods}\markdownRendererInterblockSeparator
{}\markdownRendererSectionBegin
\markdownRendererHeadingThree{2.1 Feature Engineering}\markdownRendererInterblockSeparator
{}Four groups of features were implemented, ordered by increasing representational power.\markdownRendererParagraphSeparator
{}\markdownRendererStrongEmphasis{Amino acid composition (20-d).} The normalised frequency of each canonical amino acid,\markdownRendererSoftLineBreak
{}capturing the global residue distribution. This is the simplest possible sequence\markdownRendererSoftLineBreak
{}representation and serves as a well-established baseline in protein classification literature.\markdownRendererParagraphSeparator
{}\markdownRendererStrongEmphasis{Dipeptide frequencies (400-d).} Normalised counts of all ordered amino acid pairs,\markdownRendererSoftLineBreak
{}encoding local sequential context that single-residue composition cannot capture. Combined\markdownRendererSoftLineBreak
{}with a single sequence-length scalar, this yields a 421-dimensional handcrafted composition\markdownRendererSoftLineBreak
{}vector.\markdownRendererParagraphSeparator
{}\markdownRendererStrongEmphasis{Physicochemical properties (8-d).} Molecular weight, isoelectric point, aromaticity\markdownRendererSoftLineBreak
{}(Phe + Trp + Tyr fraction), GRAVY hydropathicity, charge at pH 7.0, and Chou–Fasman\markdownRendererSoftLineBreak
{}secondary-structure propensity fractions (helix, turn, sheet). These were computed with\markdownRendererSoftLineBreak
{}BioPython's \markdownRendererCodeSpan{ProtParam} module. Unlike composition features, these capture higher-level\markdownRendererSoftLineBreak
{}biophysical properties that may distinguish enzyme families with different catalytic\markdownRendererSoftLineBreak
{}environments.\markdownRendererParagraphSeparator
{}\markdownRendererStrongEmphasis{Protein language model embeddings.} ESM-2 (Lin et al., 2023), a transformer-based PLM\markdownRendererSoftLineBreak
{}pre-trained on millions of protein sequences via masked language modelling, was used as a\markdownRendererSoftLineBreak
{}feature extractor. Two model sizes were evaluated: the 8-million-parameter variant\markdownRendererSoftLineBreak
{}(\markdownRendererCodeSpan{esm2\markdownRendererUnderscore{}t6\markdownRendererUnderscore{}8M\markdownRendererUnderscore{}UR50D}, 320-d output) and the 650-million-parameter variant\markdownRendererSoftLineBreak
{}(\markdownRendererCodeSpan{esm2\markdownRendererUnderscore{}t33\markdownRendererUnderscore{}650M\markdownRendererUnderscore{}UR50D}, 1 280-d output). For each sequence, the final-layer hidden states\markdownRendererSoftLineBreak
{}were mean-pooled across all non-special tokens to produce a single dense vector. Because the\markdownRendererSoftLineBreak
{}ESM-2 forward pass is entirely stateless—no label information is used—embeddings were safely\markdownRendererSoftLineBreak
{}pre-computed once and cached to disk. The key advantage of PLM embeddings over handcrafted\markdownRendererSoftLineBreak
{}features is that they encode evolutionary and structural information learned from millions of\markdownRendererSoftLineBreak
{}diverse protein families, capturing patterns that are impossible to design manually.\markdownRendererInterblockSeparator
{}
\markdownRendererSectionEnd \markdownRendererSectionBegin
\markdownRendererHeadingThree{2.2 Experimental Design}\markdownRendererInterblockSeparator
{}\markdownRendererStrongEmphasis{Cross-validation protocol.} All experiments used stratified five-fold cross-validation\markdownRendererSoftLineBreak
{}(seed = 42), guaranteeing that every fold preserved the original class proportions.\markdownRendererParagraphSeparator
{}\markdownRendererStrongEmphasis{Data-leakage prevention.} Three strict rules were enforced: (i) \markdownRendererCodeSpan{StandardScaler} was fit\markdownRendererSoftLineBreak
{}exclusively on the training fold and applied to both folds via \markdownRendererCodeSpan{.transform()}; (ii) SMOTE\markdownRendererSoftLineBreak
{}oversampling was applied only to the training fold; (iii) features were derived solely from\markdownRendererSoftLineBreak
{}the amino acid sequence, never from the label.\markdownRendererParagraphSeparator
{}\markdownRendererStrongEmphasis{Metrics.} Four complementary metrics were recorded for every experiment: accuracy, macro F1,\markdownRendererSoftLineBreak
{}balanced accuracy, and the Matthews Correlation Coefficient (MCC). Macro F1 was used as the\markdownRendererSoftLineBreak
{}primary model-selection criterion because it penalises poor performance on minority classes\markdownRendererSoftLineBreak
{}without being dominated by the majority class.\markdownRendererInterblockSeparator
{}
\markdownRendererSectionEnd \markdownRendererSectionBegin
\markdownRendererHeadingThree{2.3 Models}\markdownRendererInterblockSeparator
{}\markdownRendererStrongEmphasis{Baseline models.} Logistic Regression (\markdownRendererCodeSpan{class\markdownRendererUnderscore{}weight='balanced'}, multinomial, L-BFGS\markdownRendererSoftLineBreak
{}solver, \markdownRendererCodeSpan{max\markdownRendererUnderscore{}iter=1000}) and Random Forest (\markdownRendererCodeSpan{n\markdownRendererUnderscore{}estimators=300}, \markdownRendererCodeSpan{class\markdownRendererUnderscore{}weight='balanced'})\markdownRendererSoftLineBreak
{}were trained on the 429-dimensional handcrafted feature vector. A pipeline variant applied\markdownRendererSoftLineBreak
{}PCA to 50 components before the Random Forest to mitigate the curse of dimensionality imposed\markdownRendererSoftLineBreak
{}by the 400-d dipeptide features.\markdownRendererParagraphSeparator
{}\markdownRendererStrongEmphasis{Gradient-boosted trees.} XGBoost and LightGBM were trained on ESM-2 650M embeddings\markdownRendererSoftLineBreak
{}concatenated with the 8-d physicochemical features (1 288-d total). Hyperparameters for\markdownRendererSoftLineBreak
{}XGBoost were: \markdownRendererCodeSpan{max\markdownRendererUnderscore{}depth=6}, \markdownRendererCodeSpan{n\markdownRendererUnderscore{}estimators=800}, \markdownRendererCodeSpan{learning\markdownRendererUnderscore{}rate=0.1}, \markdownRendererCodeSpan{subsample=0.8},\markdownRendererSoftLineBreak
{}\markdownRendererCodeSpan{colsample\markdownRendererUnderscore{}bytree=0.9}. Class imbalance was addressed through balanced sample weights and,\markdownRendererSoftLineBreak
{}separately, through SMOTE on the training fold (five nearest neighbours, default ratio).\markdownRendererParagraphSeparator
{}\markdownRendererStrongEmphasis{End-to-end fine-tuning.} The ESM-2 8M model was fine-tuned with all transformer layers\markdownRendererSoftLineBreak
{}unfrozen. A classification head—\markdownRendererCodeSpan{Linear(320, 128) → GELU → Dropout(0.3) → Linear(128, 7)}—\markdownRendererSoftLineBreak
{}was appended after mean-pooling. Training used AdamW with differential learning rates\markdownRendererSoftLineBreak
{}(backbone: 2 × 10⁻⁵; head: 1 × 10⁻⁴), OneCycleLR scheduling, bfloat16 mixed precision,\markdownRendererSoftLineBreak
{}gradient accumulation to an effective batch size of 16, and early stopping on validation\markdownRendererSoftLineBreak
{}macro F1 with patience of 5 epochs. Per-fold class weights were computed from training-fold\markdownRendererSoftLineBreak
{}frequencies and applied to the cross-entropy loss.\markdownRendererParagraphSeparator
{}\markdownRendererStrongEmphasis{Ensemble.} A soft-vote ensemble combined the out-of-fold (OOF) probability vectors of the\markdownRendererSoftLineBreak
{}fine-tuned ESM-2 8M model and the XGBoost + SMOTE model. Weights were set proportional to\markdownRendererSoftLineBreak
{}each model's individual OOF macro F1 (w\markdownRendererUnderscore{}finetune ≈ 0.46, w\markdownRendererUnderscore{}xgboost ≈ 0.54). Per-class\markdownRendererSoftLineBreak
{}decision thresholds were optimised on the blended OOF probabilities via\markdownRendererSoftLineBreak
{}\markdownRendererCodeSpan{scipy.optimize.minimize}.\markdownRendererInterblockSeparator
{}
\markdownRendererSectionEnd \markdownRendererSectionBegin
\markdownRendererHeadingThree{2.4 Feature Ablation}\markdownRendererInterblockSeparator
{}To quantify the contribution of each feature group, XGBoost (balanced sample weights, no SMOTE)\markdownRendererSoftLineBreak
{}was trained on four feature combinations: (A) handcrafted only (429-d), (B) ESM-2 650M only\markdownRendererSoftLineBreak
{}(1 280-d), (C) ESM-2 + all handcrafted (1 709-d), and (D) ESM-2 + physicochemical only\markdownRendererSoftLineBreak
{}(1 288-d).\markdownRendererInterblockSeparator
{}
\markdownRendererSectionEnd 
\markdownRendererSectionEnd \markdownRendererSectionBegin
\markdownRendererHeadingTwo{3. Results}\markdownRendererInterblockSeparator
{}\markdownRendererSectionBegin
\markdownRendererHeadingThree{3.1 Model Comparison}\markdownRendererInterblockSeparator
{}Table 1 reports stratified five-fold cross-validation results (mean ± standard deviation).\markdownRendererParagraphSeparator
{}\markdownRendererPipe{} Model \markdownRendererPipe{} Accuracy \markdownRendererPipe{} Macro F1 \markdownRendererPipe{} Bal. Acc. \markdownRendererPipe{} MCC \markdownRendererPipe{}\markdownRendererSoftLineBreak
{}\markdownRendererPipe{}-------\markdownRendererPipe{}----------\markdownRendererPipe{}----------\markdownRendererPipe{}-----------\markdownRendererPipe{}-----\markdownRendererPipe{}\markdownRendererSoftLineBreak
{}\markdownRendererPipe{} Logistic Regression \markdownRendererPipe{} 0.457 ± 0.003 \markdownRendererPipe{} 0.191 ± 0.002 \markdownRendererPipe{} 0.300 ± 0.010 \markdownRendererPipe{} 0.171 ± 0.003 \markdownRendererPipe{}\markdownRendererSoftLineBreak
{}\markdownRendererPipe{} Random Forest \markdownRendererPipe{} 0.816 ± 0.001 \markdownRendererPipe{} 0.141 ± 0.002 \markdownRendererPipe{} 0.149 ± 0.001 \markdownRendererPipe{} 0.093 ± 0.011 \markdownRendererPipe{}\markdownRendererSoftLineBreak
{}\markdownRendererPipe{} RF + PCA(50) \markdownRendererPipe{} 0.816 ± 0.002 \markdownRendererPipe{} 0.197 ± 0.002 \markdownRendererPipe{} 0.185 ± 0.001 \markdownRendererPipe{} 0.216 ± 0.006 \markdownRendererPipe{}\markdownRendererSoftLineBreak
{}\markdownRendererPipe{} XGBoost \markdownRendererPipe{} 0.886 ± 0.003 \markdownRendererPipe{} 0.575 ± 0.013 \markdownRendererPipe{} 0.516 ± 0.011 \markdownRendererPipe{} 0.621 ± 0.010 \markdownRendererPipe{}\markdownRendererSoftLineBreak
{}\markdownRendererPipe{} LightGBM \markdownRendererPipe{} 0.879 ± 0.003 \markdownRendererPipe{} 0.515 ± 0.013 \markdownRendererPipe{} 0.467 ± 0.008 \markdownRendererPipe{} 0.605 ± 0.010 \markdownRendererPipe{}\markdownRendererSoftLineBreak
{}\markdownRendererPipe{} LightGBM + SMOTE \markdownRendererPipe{} 0.880 ± 0.002 \markdownRendererPipe{} 0.576 ± 0.014 \markdownRendererPipe{} 0.535 ± 0.010 \markdownRendererPipe{} 0.620 ± 0.006 \markdownRendererPipe{}\markdownRendererSoftLineBreak
{}\markdownRendererPipe{} XGBoost + SMOTE \markdownRendererPipe{} 0.886 ± 0.002 \markdownRendererPipe{} 0.595 ± 0.014 \markdownRendererPipe{} 0.552 ± 0.011 \markdownRendererPipe{} 0.631 ± 0.006 \markdownRendererPipe{}\markdownRendererSoftLineBreak
{}\markdownRendererPipe{} ESM-2 8M Fine-tune \markdownRendererPipe{} 0.823 ± 0.008 \markdownRendererPipe{} 0.509 ± 0.013 \markdownRendererPipe{} 0.570 ± 0.015 \markdownRendererPipe{} 0.553 ± 0.008 \markdownRendererPipe{}\markdownRendererSoftLineBreak
{}\markdownRendererPipe{} \markdownRendererStrongEmphasis{Ensemble} \markdownRendererPipe{} \markdownRendererStrongEmphasis{0.890} \markdownRendererPipe{} \markdownRendererStrongEmphasis{0.625} \markdownRendererPipe{} \markdownRendererStrongEmphasis{0.604} \markdownRendererPipe{} \markdownRendererStrongEmphasis{0.659} \markdownRendererPipe{}\markdownRendererParagraphSeparator
{}\markdownRendererEmphasis{Table 1. Cross-validation results. Baseline models use 429-d handcrafted features; advanced\markdownRendererSoftLineBreak
{}models use ESM-2 650M + physicochemical (1 288-d); fine-tuning operates on raw sequences.\markdownRendererSoftLineBreak
{}Ensemble combines fine-tuned 8M and XGBoost 650M via soft voting.}\markdownRendererParagraphSeparator
{}The progression from handcrafted baselines (macro F1 ≈ 0.19) to ESM-2-based models (0.58–0.60)\markdownRendererSoftLineBreak
{}and finally the ensemble (0.625) demonstrates the critical role of learned protein\markdownRendererSoftLineBreak
{}representations.\markdownRendererInterblockSeparator
{}
\markdownRendererSectionEnd \markdownRendererSectionBegin
\markdownRendererHeadingThree{3.2 Feature Ablation}\markdownRendererInterblockSeparator
{}\markdownRendererPipe{} Feature Set \markdownRendererPipe{} Dimensions \markdownRendererPipe{} Macro F1 \markdownRendererPipe{} Δ vs. ESM-2 only \markdownRendererPipe{}\markdownRendererSoftLineBreak
{}\markdownRendererPipe{}-------------\markdownRendererPipe{}-----------\markdownRendererPipe{}----------\markdownRendererPipe{}-------------------\markdownRendererPipe{}\markdownRendererSoftLineBreak
{}\markdownRendererPipe{} Handcrafted only \markdownRendererPipe{} 429 \markdownRendererPipe{} 0.206 \markdownRendererPipe{} −0.356 \markdownRendererPipe{}\markdownRendererSoftLineBreak
{}\markdownRendererPipe{} ESM-2 only \markdownRendererPipe{} 1 280 \markdownRendererPipe{} 0.562 \markdownRendererPipe{} — \markdownRendererPipe{}\markdownRendererSoftLineBreak
{}\markdownRendererPipe{} ESM-2 + Handcrafted \markdownRendererPipe{} 1 709 \markdownRendererPipe{} 0.552 \markdownRendererPipe{} −0.010 \markdownRendererPipe{}\markdownRendererSoftLineBreak
{}\markdownRendererPipe{} ESM-2 + Physicochemical \markdownRendererPipe{} 1 288 \markdownRendererPipe{} 0.575 \markdownRendererPipe{} +0.013 \markdownRendererPipe{}\markdownRendererParagraphSeparator
{}\markdownRendererEmphasis{Table 2. Feature ablation (XGBoost, balanced weights). Adding the full handcrafted vector\markdownRendererSoftLineBreak
{}introduces noise; the 8-d physicochemical subset provides modest complementary signal.}\markdownRendererParagraphSeparator
{}ESM-2 embeddings account for the vast majority of predictive power. The 400-d dipeptide\markdownRendererSoftLineBreak
{}features are largely redundant once PLM embeddings are available—their addition slightly\markdownRendererSoftLineBreak
{}degrades performance, likely because XGBoost overfits to the high-dimensional noise. The 8-d\markdownRendererSoftLineBreak
{}physicochemical features, particularly molecular weight, provide a small but consistent\markdownRendererSoftLineBreak
{}improvement.\markdownRendererInterblockSeparator
{}
\markdownRendererSectionEnd \markdownRendererSectionBegin
\markdownRendererHeadingThree{3.3 Confusion Matrix Analysis}\markdownRendererInterblockSeparator
{}The normalised confusion matrix of the best single model (XGBoost + SMOTE) reveals that\markdownRendererSoftLineBreak
{}class 0 (Not enzyme) achieves 93.8 \markdownRendererPercentSign{} recall, while the two smallest classes—Lyase (19.2 \markdownRendererPercentSign{}\markdownRendererSoftLineBreak
{}recall) and Isomerase (20.0 \markdownRendererPercentSign{} recall)—are most frequently misclassified as Not enzyme.\markdownRendererSoftLineBreak
{}Transferase (65.3 \markdownRendererPercentSign{} recall) and Oxidoreductase (62.9 \markdownRendererPercentSign{} recall) perform moderately. Hydrolase\markdownRendererSoftLineBreak
{}(56.5 \markdownRendererPercentSign{} recall) is sometimes confused with Transferase, consistent with the biochemical\markdownRendererSoftLineBreak
{}reality that some enzymes have overlapping catalytic mechanisms.\markdownRendererInterblockSeparator
{}
\markdownRendererSectionEnd \markdownRendererSectionBegin
\markdownRendererHeadingThree{3.4 Confidence Calibration}\markdownRendererInterblockSeparator
{}Mapping the maximum predicted probability to three confidence tiers yields the following\markdownRendererSoftLineBreak
{}calibration:\markdownRendererParagraphSeparator
{}\markdownRendererPipe{} Level \markdownRendererPipe{} Count (\markdownRendererPercentSign{}) \markdownRendererPipe{} Accuracy \markdownRendererPipe{} Mean Max Prob \markdownRendererPipe{}\markdownRendererSoftLineBreak
{}\markdownRendererPipe{}-------\markdownRendererPipe{}-----------\markdownRendererPipe{}----------\markdownRendererPipe{}---------------\markdownRendererPipe{}\markdownRendererSoftLineBreak
{}\markdownRendererPipe{} High (p ≥ 0.80) \markdownRendererPipe{} 30 448 (76.6 \markdownRendererPercentSign{}) \markdownRendererPipe{} 95.1 \markdownRendererPercentSign{} \markdownRendererPipe{} 0.959 \markdownRendererPipe{}\markdownRendererSoftLineBreak
{}\markdownRendererPipe{} Medium (0.50 ≤ p < 0.80) \markdownRendererPipe{} 6 658 (16.7 \markdownRendererPercentSign{}) \markdownRendererPipe{} 66.2 \markdownRendererPercentSign{} \markdownRendererPipe{} 0.663 \markdownRendererPipe{}\markdownRendererSoftLineBreak
{}\markdownRendererPipe{} Low (p < 0.50) \markdownRendererPipe{} 2 658 (6.7 \markdownRendererPercentSign{}) \markdownRendererPipe{} 41.2 \markdownRendererPercentSign{} \markdownRendererPipe{} 0.411 \markdownRendererPipe{}\markdownRendererParagraphSeparator
{}\markdownRendererEmphasis{Table 3. Confidence calibration on OOF predictions (XGBoost + SMOTE). The model is\markdownRendererSoftLineBreak
{}well-calibrated: high-confidence predictions are correct 95 \markdownRendererPercentSign{} of the time.}\markdownRendererInterblockSeparator
{}
\markdownRendererSectionEnd 
\markdownRendererSectionEnd \markdownRendererSectionBegin
\markdownRendererHeadingTwo{4. Interpretability}\markdownRendererInterblockSeparator
{}Four interpretability techniques were applied to the best single model (XGBoost + SMOTE).\markdownRendererParagraphSeparator
{}\markdownRendererStrongEmphasis{Feature importance.} Gini-based importance ranking shows that ESM-2 embedding dimensions\markdownRendererSoftLineBreak
{}dominate the top 20 features. The highest-ranked handcrafted feature is molecular weight\markdownRendererSoftLineBreak
{}(rank 5), confirming that ESM-2 subsumes most biochemical information.\markdownRendererParagraphSeparator
{}\markdownRendererStrongEmphasis{SHAP analysis.} TreeExplainer SHAP values (computed on a 500-sequence subsample) corroborate\markdownRendererSoftLineBreak
{}the Gini ranking: ESM-2 dimensions 112, 308, and 216 have the highest mean absolute SHAP\markdownRendererSoftLineBreak
{}values across all seven classes.\markdownRendererParagraphSeparator
{}\markdownRendererStrongEmphasis{Per-class precision, recall, and F1} (see Table 4 below) reveal that the model is\markdownRendererSoftLineBreak
{}conservative on minority classes: Isomerase precision (0.554) far exceeds its recall (0.200),\markdownRendererSoftLineBreak
{}suggesting the model correctly identifies some Isomerases but abstains from predicting this\markdownRendererSoftLineBreak
{}class in borderline cases.\markdownRendererParagraphSeparator
{}\markdownRendererPipe{} Class \markdownRendererPipe{} Precision \markdownRendererPipe{} Recall \markdownRendererPipe{} F1 \markdownRendererPipe{}\markdownRendererSoftLineBreak
{}\markdownRendererPipe{}-------\markdownRendererPipe{}-----------\markdownRendererPipe{}--------\markdownRendererPipe{}----\markdownRendererPipe{}\markdownRendererSoftLineBreak
{}\markdownRendererPipe{} Not enzyme \markdownRendererPipe{} 0.934 \markdownRendererPipe{} 0.938 \markdownRendererPipe{} 0.936 \markdownRendererPipe{}\markdownRendererSoftLineBreak
{}\markdownRendererPipe{} Oxidoreductase \markdownRendererPipe{} 0.586 \markdownRendererPipe{} 0.629 \markdownRendererPipe{} 0.607 \markdownRendererPipe{}\markdownRendererSoftLineBreak
{}\markdownRendererPipe{} Transferase \markdownRendererPipe{} 0.579 \markdownRendererPipe{} 0.653 \markdownRendererPipe{} 0.614 \markdownRendererPipe{}\markdownRendererSoftLineBreak
{}\markdownRendererPipe{} Hydrolase \markdownRendererPipe{} 0.554 \markdownRendererPipe{} 0.565 \markdownRendererPipe{} 0.559 \markdownRendererPipe{}\markdownRendererSoftLineBreak
{}\markdownRendererPipe{} Lyase \markdownRendererPipe{} 0.329 \markdownRendererPipe{} 0.192 \markdownRendererPipe{} 0.242 \markdownRendererPipe{}\markdownRendererSoftLineBreak
{}\markdownRendererPipe{} Isomerase \markdownRendererPipe{} 0.554 \markdownRendererPipe{} 0.200 \markdownRendererPipe{} 0.293 \markdownRendererPipe{}\markdownRendererSoftLineBreak
{}\markdownRendererPipe{} Ligase \markdownRendererPipe{} 0.603 \markdownRendererPipe{} 0.426 \markdownRendererPipe{} 0.499 \markdownRendererPipe{}\markdownRendererParagraphSeparator
{}\markdownRendererEmphasis{Table 4. Per-class metrics for XGBoost + SMOTE (OOF predictions).}\markdownRendererParagraphSeparator
{}\markdownRendererStrongEmphasis{High-confidence error analysis.} Of 5 312 total misclassifications, 1 499 (28.2 \markdownRendererPercentSign{}) occur at\markdownRendererSoftLineBreak
{}high confidence (p ≥ 0.80). The majority class absorbs most of these: 785 high-confidence\markdownRendererSoftLineBreak
{}errors predict Not enzyme. Lyase and Isomerase contribute disproportionately to\markdownRendererSoftLineBreak
{}high-confidence errors relative to their class sizes, indicating that the model sometimes\markdownRendererSoftLineBreak
{}assigns a confident but incorrect majority-class prediction to these rare enzyme types.\markdownRendererInterblockSeparator
{}
\markdownRendererSectionEnd \markdownRendererSectionBegin
\markdownRendererHeadingTwo{5. Discussion}\markdownRendererInterblockSeparator
{}\markdownRendererSectionBegin
\markdownRendererHeadingThree{5.1 What Worked}\markdownRendererInterblockSeparator
{}The single most impactful decision was using ESM-2 650M embeddings as features for\markdownRendererSoftLineBreak
{}gradient-boosted trees, which lifted macro F1 from 0.19 (handcrafted baseline) to 0.58—a\markdownRendererSoftLineBreak
{}threefold improvement. This confirms that pre-trained protein language models encode rich\markdownRendererSoftLineBreak
{}functional information that handcrafted sequence statistics fail to capture. The ESM-2 model,\markdownRendererSoftLineBreak
{}trained on evolutionary protein data, has implicitly learned to represent structural folds,\markdownRendererSoftLineBreak
{}binding site motifs, and conserved catalytic residues—all of which are highly relevant to\markdownRendererSoftLineBreak
{}enzyme classification but impossible to engineer from sequence alone.\markdownRendererParagraphSeparator
{}SMOTE on the training fold provided a consistent +2 percentage-point improvement in macro F1\markdownRendererSoftLineBreak
{}without harming accuracy, demonstrating that synthetic minority oversampling is beneficial\markdownRendererSoftLineBreak
{}even for PLM-based features. This result is noteworthy because SMOTE operates in the\markdownRendererSoftLineBreak
{}1 288-dimensional embedding space, where the interpolation between minority-class samples\markdownRendererSoftLineBreak
{}generates plausible synthetic enzyme representations rather than arbitrary noise.\markdownRendererParagraphSeparator
{}The soft-vote ensemble achieved the highest macro F1 (0.625) by combining two models that\markdownRendererSoftLineBreak
{}process sequences through fundamentally different pathways: XGBoost on frozen 650M embeddings\markdownRendererSoftLineBreak
{}captures static feature interactions via decision-tree splits, while the fine-tuned 8M model\markdownRendererSoftLineBreak
{}learns task-specific representations through gradient-based backpropagation. Their OOF\markdownRendererSoftLineBreak
{}prediction errors are partially uncorrelated, enabling genuine complementarity. Notably, the\markdownRendererSoftLineBreak
{}ensemble's balanced accuracy (0.604) substantially exceeds that of either constituent model,\markdownRendererSoftLineBreak
{}suggesting that the two models make different mistakes on minority-class samples.\markdownRendererInterblockSeparator
{}
\markdownRendererSectionEnd \markdownRendererSectionBegin
\markdownRendererHeadingThree{5.2 Limitations}\markdownRendererInterblockSeparator
{}\markdownRendererStrongEmphasis{Minority-class performance.} Lyase (F1 = 0.24) and Isomerase (F1 = 0.29) remain\markdownRendererSoftLineBreak
{}challenging despite all imbalance-handling strategies. With only 600 and 411 training examples\markdownRendererSoftLineBreak
{}respectively, these classes are fundamentally under-represented. The model's precision for\markdownRendererSoftLineBreak
{}Isomerase (0.554) suggests it can reliably identify some members of this class, but its low\markdownRendererSoftLineBreak
{}recall (0.200) indicates it fails to recognise the majority of Isomerases. More aggressive\markdownRendererSoftLineBreak
{}strategies—focal loss to down-weight easy majority-class examples, class-conditional data\markdownRendererSoftLineBreak
{}augmentation, or hierarchical classification (first enzyme vs. non-enzyme, then among EC\markdownRendererSoftLineBreak
{}classes)—may yield improvements.\markdownRendererParagraphSeparator
{}\markdownRendererStrongEmphasis{Fine-tuning gap.} The fine-tuned ESM-2 8M model (macro F1 = 0.51) underperforms XGBoost\markdownRendererSoftLineBreak
{}on frozen ESM-2 650M embeddings (0.60). This is attributable to two factors: (i) the 8M\markdownRendererSoftLineBreak
{}model has far less representational capacity than the 650M model (6 vs. 33 transformer\markdownRendererSoftLineBreak
{}layers), and (ii) end-to-end fine-tuning on a relatively small dataset (39 764 sequences)\markdownRendererSoftLineBreak
{}with 8 million trainable parameters risks overfitting despite regularisation (dropout,\markdownRendererSoftLineBreak
{}weight decay, early stopping). Fine-tuning the 650M model would likely close this gap but\markdownRendererSoftLineBreak
{}requires substantially more GPU memory (24 GB+).\markdownRendererParagraphSeparator
{}\markdownRendererStrongEmphasis{Handcrafted feature noise.} Adding the full 429-d handcrafted vector to ESM-2 embeddings\markdownRendererSoftLineBreak
{}slightly degraded performance (macro F1: 0.552 vs. 0.562 for ESM-2 alone). The 400-d\markdownRendererSoftLineBreak
{}dipeptide features, while informative in isolation, become redundant in the presence of PLM\markdownRendererSoftLineBreak
{}embeddings and introduce noise dimensions that XGBoost partially overfits to. The 8-d\markdownRendererSoftLineBreak
{}physicochemical subset, by contrast, provides complementary signal (macro F1: 0.575),\markdownRendererSoftLineBreak
{}confirming that carefully selected biophysical descriptors retain value even alongside\markdownRendererSoftLineBreak
{}powerful PLM representations.\markdownRendererInterblockSeparator
{}
\markdownRendererSectionEnd \markdownRendererSectionBegin
\markdownRendererHeadingThree{5.3 Error Analysis}\markdownRendererInterblockSeparator
{}The confusion matrix reveals biologically interpretable patterns. Hydrolase is sometimes\markdownRendererSoftLineBreak
{}confused with Transferase; this is plausible because both catalyse bond cleavage, differing\markdownRendererSoftLineBreak
{}primarily in whether the cleaved group is transferred to water (hydrolysis) or to another\markdownRendererSoftLineBreak
{}acceptor molecule (transfer). Lyase and Isomerase errors predominantly default to Not enzyme,\markdownRendererSoftLineBreak
{}reflecting the model's conservative behaviour under class imbalance—when uncertain, it\markdownRendererSoftLineBreak
{}predicts the overwhelming majority class.\markdownRendererParagraphSeparator
{}Of 5 312 total misclassifications, 1 499 (28.2 \markdownRendererPercentSign{}) occur at high confidence (p ≥ 0.80).\markdownRendererSoftLineBreak
{}These high-confidence errors are concentrated on class 0 predictions (785 errors), suggesting\markdownRendererSoftLineBreak
{}that the model occasionally assigns very high probability to the majority class even for\markdownRendererSoftLineBreak
{}genuine enzymes. This is consistent with the well-documented calibration challenges of\markdownRendererSoftLineBreak
{}tree-based models on imbalanced data, where the learned decision boundaries are biased toward\markdownRendererSoftLineBreak
{}the majority class.\markdownRendererInterblockSeparator
{}
\markdownRendererSectionEnd 
\markdownRendererSectionEnd \markdownRendererSectionBegin
\markdownRendererHeadingTwo{6. Conclusion}\markdownRendererInterblockSeparator
{}This project demonstrates that protein language model embeddings, combined with\markdownRendererSoftLineBreak
{}gradient-boosted trees and careful class-imbalance handling, achieve strong performance\markdownRendererSoftLineBreak
{}on seven-class enzyme classification from sequence alone. The systematic progression from a\markdownRendererSoftLineBreak
{}macro F1 of 0.19 (handcrafted baseline) to 0.625 (ensemble) confirms the dominance of\markdownRendererSoftLineBreak
{}learned protein representations over traditional feature engineering. The ensemble of a\markdownRendererSoftLineBreak
{}fine-tuned ESM-2 8M model and XGBoost on frozen 650M embeddings exploits complementary model\markdownRendererSoftLineBreak
{}strengths and achieves the best overall performance across all four evaluation metrics.\markdownRendererParagraphSeparator
{}The ablation study demonstrates that ESM-2 embeddings are the dominant source of predictive\markdownRendererSoftLineBreak
{}power, while a small set of physicochemical features provides complementary biophysical\markdownRendererSoftLineBreak
{}signal. SMOTE oversampling on training folds consistently improves minority-class recall\markdownRendererSoftLineBreak
{}without degrading overall accuracy. Confidence calibration confirms that the model is\markdownRendererSoftLineBreak
{}well-calibrated, with high-confidence predictions achieving 95 \markdownRendererPercentSign{} accuracy.\markdownRendererParagraphSeparator
{}Key directions for future work include: (i) fine-tuning the larger 650M model, which would\markdownRendererSoftLineBreak
{}likely close the fine-tuning gap and further improve ensemble performance; (ii) applying focal\markdownRendererSoftLineBreak
{}loss or cost-sensitive learning to improve performance on the most difficult minority classes\markdownRendererSoftLineBreak
{}(Lyase, Isomerase); and (iii) hierarchical classification—first distinguishing enzyme from\markdownRendererSoftLineBreak
{}non-enzyme, then classifying among the six EC classes—which could exploit the natural two-level\markdownRendererSoftLineBreak
{}structure of the label space.\markdownRendererParagraphSeparator
{}The complete codebase, including the blind prediction pipeline, is fully reproducible with a\markdownRendererSoftLineBreak
{}single seed (42) and ready for deployment on unseen test data.\markdownRendererInterblockSeparator
{}
\markdownRendererSectionEnd \markdownRendererSectionBegin
\markdownRendererHeadingTwo{7. Appendix}\markdownRendererInterblockSeparator
{}\markdownRendererSectionBegin
\markdownRendererHeadingThree{A. Blind Challenge Prediction Command}\markdownRendererInterblockSeparator
{}\markdownRendererInputFencedCode{./_markdown_REPORT/3a741ee7b03893af3e3e1fd0bba16bb0.verbatim}{}{}\markdownRendererInterblockSeparator
{}Output format (one line per sequence, no header):\markdownRendererInterblockSeparator
{}\markdownRendererInputFencedCode{./_markdown_REPORT/36f5def0d331205dbc3d822901c35d1a.verbatim}{}{}\markdownRendererInterblockSeparator
{}
\markdownRendererSectionEnd \markdownRendererSectionBegin
\markdownRendererHeadingThree{B. Software Dependencies}\markdownRendererInterblockSeparator
{}\markdownRendererInputFencedCode{./_markdown_REPORT/f288ddd041d473901ba813b9f60473f0.verbatim}{}{}\markdownRendererInterblockSeparator
{}
\markdownRendererSectionEnd \markdownRendererSectionBegin
\markdownRendererHeadingThree{C. Reproducibility}\markdownRendererInterblockSeparator
{}All random seeds fixed to 42 (\markdownRendererCodeSpan{random}, \markdownRendererCodeSpan{numpy}, \markdownRendererCodeSpan{torch}). Stratified five-fold CV with\markdownRendererSoftLineBreak
{}\markdownRendererCodeSpan{shuffle=True, random\markdownRendererUnderscore{}state=42}. Results stored as JSON for exact figure reproduction without\markdownRendererSoftLineBreak
{}retraining.
\markdownRendererSectionEnd 
\markdownRendererSectionEnd 
\markdownRendererSectionEnd \markdownRendererDocumentEnd